%% LyX 2.4.4 created this file.  For more info, see https://www.lyx.org/.
%% Do not edit unless you really know what you are doing.
\documentclass[english]{article}
\usepackage[T1]{fontenc}
\usepackage[utf8]{inputenc}
\setcounter{secnumdepth}{-2}
\setlength{\parindent}{0bp}
\usepackage{amsmath}
\usepackage{amsthm}
\usepackage{amssymb}

\makeatletter
%%%%%%%%%%%%%%%%%%%%%%%%%%%%%% Textclass specific LaTeX commands.
\theoremstyle{plain}
\newtheorem*{thm*}{\protect\theoremname}

\@ifundefined{date}{}{\date{}}
%%%%%%%%%%%%%%%%%%%%%%%%%%%%%% User specified LaTeX commands.
\usepackage{graphicx}
\usepackage{geometry}
\newcommand{\limi}{\underset{n\rightarrow\infty}{\lim}}
\newcommand{\limxz}{\underset{x\rightarrow x_{0}}{\lim}}
\newcommand{\limfxz}{\underset{x\rightarrow x_{0}}{\lim}f(x)}
\newcommand{\limfxm}{\underset{x\rightarrow x_{0}^{-}}{\lim}f(x)}
\newcommand{\limfxp}{\underset{x\rightarrow x_{0}^{+}}{\lim}f(x)}
\newcommand{\limxm}{\underset{x\rightarrow x_{0}^{-}}{\lim}}
\newcommand{\limxp}{\underset{x\rightarrow x_{0}^{+}}{\lim}}
\newcommand{\sqrm}{M_{m\times m}(\mathbb{F})}
\newcommand{\seqa}{(a_{n})_{n=1}^{\infty}}
\newcommand{\seqx}{(x_{n})_{n=1}^{\infty}}
\newcommand{\funcdr}{f:D\rightarrow\mathbb{R}}
\newcommand{\der}{\limxz\frac{f(x)-f(x_{0})}{x-x_{0}}}
\usepackage[all]{pdfprivacy}

\makeatother

\usepackage{babel}
\providecommand{\theoremname}{Theorem}

\begin{document}
\title{{\Large Fermat's Theorem\vspace{-2em}}}
\maketitle
\begin{thm*}
Let $f:\left(a,b\right)\rightarrow\mathbb{R}$ be a function from
$\left(a,b\right)$ to $\mathbb{R}$ and suppose that $x_{0}\in\left(a,b\right)$
is a point where $f$ has a local extremum. If $f$ is differentiable
at $x_{0}$, then $f'(x_{0})=0$.
\end{thm*}
\medskip{}

\rule[0.5ex]{1\columnwidth}{1pt}

\medskip{}

\begin{proof}
Since $x_{0}$ is a local extremum point of $f$, it's either a local
maximum point or a local minimum point. Let us assume $x_{0}$ is
a local maximum point. Then since it's also in particular an interior
point, 
\[
\exists\delta>0\quad s.t.\quad\forall x\in\left(x_{0}-\delta,x_{0}+\delta\right)\subseteq\left(a,b\right),f(x)\le f(x_{0})
\]

Since $f$ is differentiable at $x_{0}$ it's in particular differentiable
from the right and from the left while maintaining 
\[
f'_{-}(x_{0})=f'(x_{0})=f'_{+}(x_{0})
\]

We shall examine each one-sided derivative of $f$ at $x_{0}$ in
turn.

For every $x\in\left(x_{0}-\delta,x_{0}\right)$ we get $f(x)\le f(x_{0})$
and $x<x_{0}$, which gives 
\[
\forall x\in\left(x_{0}-\delta,x_{0}\right),\frac{f(x)-f(x_{0})}{x-x_{0}}\ge0
\]

Let us recall that limits weakly preserve inequalities. Hence

\[
f'_{-}\left(x_{0}\right)=\underset{x\rightarrow x_{0}^{-}}{\lim}\frac{f(x)-f(x_{0})}{x-x_{0}}\ge0
\]
For every $x\in\left(x_{0},x_{0}+\delta\right)$ we get $f(x)\le f(x_{0})$
and $x>x_{0}$, which gives 

\[
\forall x\in\left(x_{0},x_{0}+\delta\right),\frac{f(x)-f(x_{0})}{x-x_{0}}\le0
\]

And again since limits weakly preserve inequalities,

\[
f'_{+}\left(x_{0}\right)=\underset{x\rightarrow x_{0}^{+}}{\lim}\frac{f(x)-f(x_{0})}{x-x_{0}}\le0
\]
This gives us 
\[
0\le f'(x_{0})\le0
\]
Hence 
\[
f'(x_{0})=0
\]
As required. 

Now let us assume $x_{0}$ is a local minimum of $f$.

Let us recall that $c$ is a local minimum of $f$ if and only if
$c$ is a local maximum of $-f$.

Let us also recall that by the differentiable functions arithmetic
theorem, $f$ is differentiable at point $c$ and $f'(c)=\alpha$
if and only if $-f$ is differentiable at point $c$ and $\left(-f\right)'(c)=-\alpha$.

Let us notice that $\left(-f\right)$ fulfills the requirements for
the case we just proved, hence the derivative of $-f$ at point $x_{0}$
is zero, which in turn means $f'(x_{0})=0$, as required.
\end{proof}

\end{document}
